\section{Deterministic growing likelihood variances}
\label{sec:raw-growing-temperature}
The obstruction persists for slowly growing, predetermined likelihood
variances. The following statement allows a different adverse law at
each sample size, as does the main theorem. The temperature sequence
must be fixed before observing the data.

\begin{corollary}[A growing-temperature obstruction]
Fix $p>2$, $T_0>0$, and
\[
 \frac12<r<\frac p{p+1}.
\]
Let $(T_n)$ be any deterministic sequence satisfying
$T_0\le T_n\le n^{r-1/2}$ for all sufficiently large $n$.
For the same fixed one-neuron Student prior and the raw Gaussian
likelihood at temperature $T_n$, there are bounded noiseless laws
$P_n$ and constants $c_*,q_*>0$, depending only on $p,r,T_0$, such that
\[
 \begin{split}
 P_n^{\otimes n}\left\{R_{P_n}(\bar f_{S_n})
       \ge c_*\frac{n^{2r-1}}{(T_n+1)^2}\right\}&\ge q_*,\\
 R_{P_n}(0)&\le n^{r-1}.
 \end{split}
\]
The same lower bound, multiplied by $q_*$, holds in expectation.
In particular, the risk diverges for every polylogarithmically
bounded temperature sequence, whereas the zero-comparator risk
vanishes. For $p>3$, one may choose $r>3/4$; then the risk still
diverges when $T_n=n^{1/4}$.
\end{corollary}

\begin{proof}
Use the same three-point construction and balanced-count event, but
now set
\[
 b=n^{(r-1)/2},\qquad
 R=\exp\left\{\frac{n^r}{100p(T_n+1)}\right\},\qquad a=b/R.
\]
All prior density bounds in the preceding proof are independent of
temperature. Integrating the likelihood in the saturated region
now gives $C\sqrt{T_n/n}$, so that contribution is bounded by
$C a^p\sqrt{T_n/n}$. Integrating the two regular prediction
coordinates gives $CT_n/n$. The appropriate evidence scale is
\[
 Z_{*,n}:=\frac{T_n}{n}a^p b^{-p-1}
        =\frac{T_nR^{-p}}{bn}.
\]
The ratio of the saturated bound to this scale is at most
\[
 C b^{p+1}\sqrt{n/T_n}
 \le\frac C{\sqrt{T_0}}
       n^{[(r-1)(p+1)+1]/2}=o(1).
\]
On the bad-fit complement the likelihood is at most
$\exp\{-n^r/(12T_n)\}$.
Its exponential rate exceeds $p\log R=n^r/[100(T_n+1)]$ by
at least $c n^r/(T_n+1)$. Since $T_n\le n^{r-1/2}$, this last
quantity is at least $c'\sqrt n$ for large $n$. The polynomial
factors in $Z_{*,n}$ cannot offset this gap. Hence
$Z_{S_n}\le C Z_{*,n}$ uniformly over the allowed temperatures.

For the numerator use prediction-coordinate radius
$t_n=\sqrt{T_n}/(10\sqrt n)$ instead of $1/(10\sqrt n)$.
Its ratio to $b$ is
\[
 \frac{t_n}b=\frac1{10}\sqrt{T_n/n^r}
                 \le\frac1{10}n^{-1/4},
\]
so the rectangle still lies in the good-fit region. Its likelihood
is at least $e^{-1/100}$ and its area is $T_n/(25n)$.
The same density and Jacobian lower bounds, integrated over
$1\le v\le R/100$, therefore give numerator at least
$c Z_{*,n}\log R$. The negative contribution is still bounded by
$\mathbb E|v|e^{-n^r/(12T_n)}=o(Z_{*,n}\log R)$.
Consequently $\bar f_{S_n}(1)\ge c\log R$, which gives the stated
rare-input risk. All constants are uniform in the prescribed
temperature sequence within the displayed range.
\end{proof}

This corollary does not cover data-dependent temperature selection.
It also does not assert failure for every possible deterministic
temperature growth rate.
